% Experimental design -- budget 800 words (body text).
% Sources: docs/decisions.md D54, D56, D57, D60-D66, D68-D71, D76; docs/methods.md sections 6.1, 8, 9;
% src/inject.py (Table: DUR_H, _pick_zones, _place), src/benchmark.py (detection_metrics, step_auc, attribution_scores),
% src/label_model.py (ece), src/downstream.py (mase_scale, dm_test, dm_table); results/benchmark_v1/summary.json and
% dev_selection.json (8,526 events; 48 settings; development seeds 100-104).
\section{Experimental design}\label{sec:experiments}

Detection and attribution are measured on injected events, and the value of the catalogue downstream.

\subsection{Injection benchmark}\label{sec:injection}

Events are written into the real southern series inside \emph{clean zone-time}: zone-steps more than 53 steps from any rule state, after a 35-day burn-in and away from clock-change days. Table~\ref{tab:injection} lists the nine types; two untyped ones probe the detector where no rule fires. Events are placed in random type order, without overlap, until they cover 1, 3 or 7\% of the clean zone-steps.

The predictors are not retrained: cross-fitted on the real series, they predict every injected step from the other half of the real data. The rule engine, label model and end model are re-run on each injected dataset as on the real data. Every design choice was made on five development datasets (seeds 100--104). The 30 test datasets (three rates $\times$ seeds 0--9; 8,526 events) were used for no choice. Results refer to the test period, July 2023 to May 2024; the earlier period serves only to tune $\beta$.

\begin{table}[tbp]
\centering
\caption{Injected event types. Durations are log-uniform in the range, except G4 (uniform). $\xh$: expected value (Eq.~\eqref{eq:expected}); $U(a,b)$: uniform draw. Artefact signatures are copied from real episodes; grid-event signatures follow the hypothesised physics of Section~\ref{sec:typology}. Each event blocks its zones for $\pm 12$~h, and a multi-city event blocks the whole system for $\pm 3$~h. Clock errors (A4) and load transfers (G5) are not injected: their rules are calendar-gated or act on whole days.}
\label{tab:injection}
\footnotesize
\setlength{\tabcolsep}{4pt}
\begin{tabularx}{\textwidth}{@{}l L{3.6cm} l Y@{}}
\toprule
Type & Zones & Duration & Injected signature \\
\midrule
A1 & all zones of 2--3 cities & 0.5--24~h & exact zeros \\
A2 & 2 or more zones of one city (50\%), one city (25\%) or 2--3 cities (25\%) & 1--48~h & each zone held at its last value or interpolated by a straight line (50/50); a whole-city fill is followed by an A1 zero fill with probability 0.5 \\
A3 & one zone & 1--6~h & hold or straight-line fill \\
\midrule
G1 & all zones of 2 cities & 1--18~h & noise floor $\xh \cdot U(0, 0.03)$ with 0--30\% exact zeros; linear ramp of 20--90~min; overshoot $\omega \sim U(1.1, 1.5)$ decaying over 1--2~h \\
G2 & all zones of one city & 20~min--6~h & as G1 \\
G3 & one zone, any hour & 1--48~h & as G1 \\
G4 & one zone, start 06:00--11:59, 70\% at weekends & 1--6~h & noise floor as G1; ramp of 10--30~min; no overshoot \\
\midrule
U\_SHOCK & all zones of 2--3 cities & 0.5--4~h & scaled by $U(0.5, 0.8)$; no rule fires \\
U\_PARTIAL & one zone & 1--8~h & scaled by $U(0.10, 0.45)$, above the dead threshold \\
\bottomrule
\end{tabularx}
\end{table}

\subsection{Methods compared}\label{sec:compared}

Detection compares the rule engine alone, the main detector alone and the full method (rule items plus detector-only candidates) with three baselines: STL \cite{cleveland1990stl,seabold2010statsmodels} with a daily period, scored by the robust $z$-score of its remainder (Eq.~\eqref{eq:score}), and Isolation Forest \cite{liu2008iforest} and LOF \cite{breunig2000lof} on five per-step features of $\rho$. Ablations change one component of the main detector: the learned graph (removed), the predictor (LightGBM), the scales (raw band only), the hierarchy (flat zone score) or the event formation (untuned).

Every detector, baselines included, receives its own event formation: the best of the same 48 settings of $(w, q, g_{\max})$ by mean event F1 on the development datasets. Ablations inherit the main detector's setting; the no-graph ablation is also reported at its own. The main detector was chosen among four predictors by a criterion fixed beforehand: the development F1 of the rules plus its candidates. Baselines that were not run are discussed in Section~\ref{sec:discussion}.

Attribution compares the rules' majority vote, the label model, the end model and the fusion on the typed injected events that the full method detects.

\subsection{Metrics}\label{sec:metrics}

Detection is counted per event, without the point adjustment that inflates scores \cite{kim2022rigorous}. With $E$ the injected and $\hat{E}$ the predicted events,
\begin{equation}
\begin{gathered}
P = \frac{|\hat{E}^{*}|}{|\hat{E}|}, \qquad R = \frac{|E^{*}|}{|E|}, \qquad F_1 = \frac{2PR}{P + R}, \\
\hat{E}^{*} = \{\hat{e} \in \hat{E} : \hat{e} \cap E \neq \emptyset\}, \qquad E^{*} = \{e \in E : e \cap \hat{E} \neq \emptyset\},
\end{gathered}
\label{eq:eventf1}
\end{equation}
where $\hat{e} \cap E \neq \emptyset$ means that $\hat{e}$ overlaps an injected span in time, its restoration tail included; the zone-strict variant also requires a shared zone. Predictions that touch neither clean zone-time nor an injected span are not scored. A second view also ignores each method's own alarms on the uninjected series, as rule states already are. Step-wise AUC-PR and recall per type are also reported.

Attribution is scored by macro-F1 over the seven typed classes and by the expected calibration error \cite{guo2017calibration} of the arg-max probability over ten equal-width bins $B_b$:
\begin{equation}
\mathrm{ECE} = \sum_{b=1}^{10} \frac{|B_b|}{n} \left| \mathrm{acc}(B_b) - \mathrm{conf}(B_b) \right|.
\label{eq:ece}
\end{equation}

Forecasts $\hat{y}_{z,t}$ are scored on a set $S$ of zone-steps by the mean absolute scaled error \cite{hyndman2006mase},
\begin{equation}
\mathrm{MASE} = \frac{1}{|S|} \sum_{(z,t) \in S} a_{z,t}, \qquad a_{z,t} = \frac{|x_{z,t} - \hat{y}_{z,t}|}{s_z},
\label{eq:mase}
\end{equation}
where $s_z$ is the mean absolute one-day change of zone $z$ on normal training steps. Two strategies are compared by a Diebold--Mariano test \cite{diebold1995dm} on the loss summed over zones, $L_t = \sum_z a_{z,t}$:
\begin{equation}
d_t = L^{(1)}_t - L^{(2)}_t, \qquad \mathrm{DM} = \frac{\bar{d}}{\sqrt{\hat{\Omega}/n}}, \qquad
\hat{\Omega} = \hat{\gamma}_0 + 2 \sum_{k=1}^{h} \Big(1 - \frac{k}{h+1}\Big) \hat{\gamma}_k,
\label{eq:dm}
\end{equation}
where $\hat{\gamma}_k$ is the lag-$k$ autocovariance of $d_t$, with the Newey--West bandwidth \cite{newey1987hac} equal to the horizon $h$ because direct forecasts overlap. Outage counts are compared by the ratio of counted to true outage zone-hours and its log distance:
\begin{equation}
r = H / H^{\mathrm{true}}, \qquad \delta = |\log r|.
\label{eq:outage}
\end{equation}

\subsection{Downstream test}\label{sec:downstream}

On the southern system, the cleaning strategies S0--S3 (Section~\ref{sec:cleaning}) feed two forecasters, 1~h and 1~day ahead: a global LightGBM direct forecaster that also reads the event indicator, trained up to 31 December 2023, and a model-free reference (persistence at 1~h, seasonal naive at 1~day). Forecasts are scored from 1 January to 24 May 2024 against the recorded values, on zone-steps that every strategy treats as normal, and separately on those within two days after an anomalous state. Forecasting uses the real data and four injected test datasets; outage counts use all 30. Naive counting takes every flagged dead state and typed counting only those inside grid-event items; the truth is the injected G1--G4 spans.

\subsection{Transfer test}\label{sec:transfer}

Tetouan tests the method on a new system with nothing re-tuned. The rules, the label-model parameters (applied by one E-step), the architectures, the training settings and the event formation are frozen; the expected value, the score scales, the predictor weights and $\gamma_q$ (at the frozen percentile) are re-estimated from Tetouan's unlabelled series. With one city, G1 is not injected and multi-city types become city-wide. Thirty fresh injected datasets (three rates $\times$ 10 seeds) are scored over 2017.

\subsection{Statistics and reproducibility}\label{sec:statistics}

Results are mean $\pm$ standard deviation over the 10 seeds of each rate. Methods are compared by exact paired two-sided Wilcoxon signed-rank tests \cite{wilcoxon1945}: the smallest attainable $p$ is 0.002 over 10 seeds and $1.9 \times 10^{-9}$ over 30 datasets. The code uses Python, PyTorch \cite{paszke2019pytorch}, LightGBM 4.6.0 \cite{ke2017lightgbm} and scikit-learn \cite{pedregosa2011sklearn}; every injected dataset is regenerated from its rate and seed, and versions, settings and run times are given in the supplementary material.
